\documentclass[12pt]{article}
\usepackage[utf8]{inputenc}
\usepackage[T1]{fontenc}
\usepackage{hyperref}
\title{Sicana Language Resources for K-12 Curriculum Development}
\author{Compiled by Claude Sonnet 5.5 (claude-sonnet-5-5) for Andr\'as Kornai}
\date{}
\begin{document}
\maketitle

\section{Introduction}
The search found no resources that match the name ``Sicana'' or the code sxc directly. The nearest match is Sicanian, the poorly attested ancient language of the Sicani in western and central Sicily. This identification is an inference from the name and the ``unclassifiable'' label. It is not confirmed by any source found. Search results describe Sicanian as known only from a few short inscriptions written in the Greek alphabet. Apart from names, these inscriptions are largely untranslated. Scholars dispute whether Sicanian was a distinct language at all, and whether it was still spoken in the classical period. There are no native speakers, no community, and no teaching materials. The search also returned pages on Sikiana (a Cariban language) and Sikka. These are different languages and are not used here.

\section{Primary Educational Resources}
No primers, textbooks, or school programmes exist for this language. None were found in search results.

\section{Linguistic and Reference Materials}
\begin{itemize}
\item Encyclopedia article on the Sicani people. It gives background on the people and the language question. \url{https://wikipedia.com/wiki/Sicani}
\item ``The non-Greek languages of Sicily'', a page from the PRIN Italia antica project. It surveys the non-Greek languages of Sicily, which would include Sicanian. I did not read it in full. \url{https://www.prin-italia-antica.unifi.it/p152.html}
\item Front matter of the Cambridge volume \emph{Language and Linguistic Contact in Ancient Sicily}, edited by Olga Tribulato. It is a scholarly entry point to the contact situation. \url{https://assets.cambridge.org/97811070/29316/frontmatter/9781107029316_frontmatter.pdf}
\end{itemize}

\section{Religious and Community Resources}
No religious texts or community organizations were found. No speaker community exists.

\section{Digital and Audio Resources}
No audio, video, corpus, or software resources were found.

\section{Curriculum Development Context}
A curriculum taught in Sicanian is not feasible. The few inscriptions and the personal and place names are too limited to support language teaching. The reference works above could support a short unit on the ancient languages and peoples of Sicily within a history or classics course, such as the Sicani, Elymians, and Sicels and their contact with Greek. Any such unit should state plainly that Sicanian is unclassified and almost undocumented. The identification of the sxc entry with ancient Sicanian should be checked against the ISO 639-3 registry before use.
\end{document}